\section*{Chapter \ref{ch:iv2:your-strategy-and-your-project} --- Your Strategy and Your Project}
\markboth{\protect\omTocMark{}{Answers: Your Strategy and Your Project}}{\protect\omTocMark{}{Answers: Your Strategy and Your Project}}

\begin{solution}{iq:iv2:your-strategy-and-your-project:1}
Say what the problem was, what you did, what came of it, in numbers, and what you would question in it. Then stop, in under two minutes, and let them pick where to go.

To build your own, write one sentence each on the problem, the approach, your own part, the result, and the main weakness. Time it out loud. Check each number: you must be able to derive it at a board, and to say how you tested it.

\textbf{Illustration, with your own facts in the brackets.} ``I wanted to know [question], because [why it mattered]. I used [data] and [method]. My part was [what you did]. The result was [number] against a baseline of [number]. I trust it less for [weakness], and I'd test that with [next step].''

\iqwrong{Telling it in the order you did it. It takes ten minutes, and the result arrives last, if at all.}
\iqfollow{Why that method rather than something simpler?}{Say what the simple one gave. If you never ran it, say so, and how you'd compare them.}
\end{solution}

\begin{solution}{iq:iv2:your-strategy-and-your-project:2}
About 2 points of alpha, at most. The market part was $1.4 \times 20\% = 28\%$, which leaves $30 - 28 = 2\%$.

And a single year's 2\% is well inside the noise. So the good answer says that this year doesn't show skill, and talks about the process instead.

\iqwrong{Comparing 30 with 20 and claiming 10 points. If you held 1.4 times the market, you've earned the market's return on 1.4 times the money.}
\end{solution}

\begin{solution}{iq:iv2:your-strategy-and-your-project:3}
The standard error is about 1.6 points, so 60\% is about 6.3 standard errors above a coin. Yes, it's better than a coin on that test. The next question is whether the test set was clean.

With $p = 0.5$ under the null, $\sqrt{0.25/1000} = 1.58\%$, and $10/1.58 = 6.3$.

So what comes next: how many models did you try on that test set, how were the cases split, and is 60\% worth anything after costs? A significant result can still be useless.

\iqwrong{Stopping at the $z$-score. If the same test set chose the model, or if nearby rows leaked into training, the 6.3 means nothing.}
\end{solution}

\begin{solution}{iq:iv2:your-strategy-and-your-project:4}
Name one specific decision, say what you now know that you didn't, and say what it would have changed. Choose a real one that costs you something: an interviewer doesn't believe ``nothing'' or ``more time''.

To build your own, list the decisions that took the most time or that you reversed. Pick one you can explain with its effect, ideally with a number.

\textbf{Illustration, with your own facts in the brackets.} ``I'd have [built the baseline or the data checks] first. I started with [complex method], and I found [problem] only at [late stage]. With [simple check] at the start, I'd have saved [time] and found [result] earlier.''

\iqwrong{A disguised boast, such as ``I'd have aimed even higher''. It tells the interviewer that you don't look for your own errors.}
\end{solution}

\begin{solution}{iq:iv2:your-strategy-and-your-project:5}
Name it, and say why: the sample was short, a choice was made after seeing the data, or it depends on one assumption. Then say what test would settle it. Doubting one of your own results is the strongest evidence of your judgement, and it makes the others easier to believe.

To build your own, go through each number you plan to mention, and ask which one you'd attack if a rival had written it.

\textbf{Illustration, with your own facts in the brackets.} ``The [result] is the weakest. It rests on [short sample or one assumption], and I chose [parameter] after seeing [what]. I would rerun it with [fresh data or a fixed choice] before I'd rely on it.''

\iqwrong{Saying that you believe them all equally. Nobody with experience believes that of their own work, and the next question will find the weak one for you.}
\end{solution}

\begin{solution}{iq:iv2:your-strategy-and-your-project:6}
Yes, it could be. A true Sharpe of zero gives a record this good about 3.3\% of the time, so about 7 of the 200 would show it with no skill at all.

The $t$-statistic is $1.5\sqrt{1.5} = 1.84$, and $\Phi(-1.84) = 3.3\%$. So $200 \times 3.3\% = 6.6$.

So the record gets you the conversation, not the verdict. What helps is to show the process: why you trade what you trade, and how your actual trades compare with your plan. Also tell them how many other things you tried that you dropped.

\iqwrong{Treating 18 good months as proof. Interviewers see the candidates who survived, and so do you.}
\end{solution}

\begin{solution}{iq:iv2:your-strategy-and-your-project:7}
Give the result in one sentence that a smart non-specialist follows, then the method in three. Then be honest about use: name the skill that carries over, even if the topic doesn't.

To build your own, write the thesis at three lengths: one sentence, one minute and ten minutes. Prepare for a mathematician's question on the core argument, and for the question ``so what?''.

\textbf{Illustration, with your own facts in the brackets.} ``I showed that [result], which matters because [plain reason]. The topic itself won't make money here. What carries over is [estimating from noisy data, or proving something works under limited assumptions], which I'd use on [specific problem].''

\iqwrong{Forcing a link to trading that isn't there. An interviewer knows, and it costs you more than saying that the topic has no use here.}
\iqfollow{What's the hardest step in the proof, or the main step of the estimate?}{Be ready to do it at the board, without notes. That's the check that it's yours.}
\end{solution}

\begin{solution}{iq:iv2:your-strategy-and-your-project:8}
Name a type of counterparty, and why they would trade at this price for a reason other than value: a forced seller, an index that has to rebalance, a hedger paying for insurance, or a slower reader of the news. If you can't name them, the likeliest loser is you.

To build your own, finish the sentence ``[this group] trades this because [rule or constraint], so the price is off by about [size], until [event].'' Then ask why a larger fund hasn't already taken the trade.

\textbf{Illustration, with your own facts in the brackets.} ``The other side is [group], who must [constraint] by [date] regardless of price. That leaves the price about [size] from [fair value]. It's too small for [large players] because [reason]. If [group]'s constraint goes away, so does my reason for the trade.''

\iqwrong{``The market is wrong.'' That says nothing about who's wrong or why they'd stay wrong. An interviewer hears a view, but no edge.}
\end{solution}

\begin{solution}{iq:iv2:your-strategy-and-your-project:9}
Stop and check it. Say ``Let me look'', reread the step, and if they're right, say so at once and without fuss. Then rebuild from that point: what falls, what stays, and what you would now conclude.

\begin{enumerate}
\item Don't defend it while you're checking it.
\item Say what you got wrong in one clear sentence.
\item Rework the part it touches, out loud, and give the new conclusion.
\item If you still disagree after checking, say why, with a small example they can verify.
\end{enumerate}

To build your own habit, practise with someone who interrupts you with real errors. Being corrected is often planted: they want to see this minute more than the answer.

\iqwrong{Arguing for several minutes, or collapsing and giving up on the question. Either one tells them how you'd behave with a losing position, which is the thing they're trying to learn.}
\iqfollow{What if you think that they're the one who's wrong?}{Say so politely and test it with a small numerical case. If they're testing you, standing firm on a checked step is what they're looking for.}
\end{solution}

\begin{solution}{iq:iv2:your-strategy-and-your-project:10}
Decline politely, and describe what you can in general terms. You don't share a former employer's confidential work, and telling them shows that you'd do the same with theirs.

To build your own, decide before the interview what's public or generic: the kind of problem, the type of data, and how you tested and ran it. Check what your contract says, and don't guess.

\textbf{Illustration, with your own facts in the brackets.} ``I can't go into the specifics, as they belong to [employer]. In general terms, I worked on [public category of signal], over [broad horizon], and I was responsible for [your part]. I can show you how I'd approach it from scratch using public methods.''

\iqwrong{Hinting, with a wink, at the signals. A firm that hears this assumes that you'd do the same to them, and it may also be taking on legal risk by listening.}
\iqfollow{We just want a flavour. Surely you can tell us the horizon and the universe?}{Only what you've checked is public or generic. Repeat your offer to work the problem from scratch.}
\end{solution}

\begin{solution}{iq:iv2:your-strategy-and-your-project:11}
Say exactly what you did, and name what others did. The word is ``I'' for your part, and ``we'' or a role for theirs. If you can't reproduce the number alone, say which part you'd need help with.

To build your own, write down your part of each result on your CV. If the honest part is small, then lead with what you learned, and how you'd do it again.

\textbf{Illustration, with your own facts in the brackets.} ``The team improved [metric] by [number]. I built [component], which accounted for [what you can honestly claim, or say that you can't separate it]. [Role] did [other part]. I could rebuild [component] today, but I'd need to relearn [other part].''

\iqwrong{``We'' throughout, or ``I'' throughout. The first leaves them guessing what you did, and the second falls apart at the first detail. Interviewers check by going two levels down.}
\end{solution}

\begin{solution}{iq:iv2:your-strategy-and-your-project:12}
The expected return is 2.8\%. The trade stays positive while your 60\% is above about 44.4\%. Kelly gives 3.5 times leverage. I wouldn't use that: I'd use well under half of it, given how uncertain the 60\% is.

$0.6 \times 10 - 0.4 \times 8 = 2.8$. Break-even is $8/18 = 44.4\%$.

For a two-outcome trade with win $w$ and loss $l$, $f = p/l - (1 - p)/w = 7.5 - 4 = 3.5$.

At 3.5 times, the 8\% loss costs you 28\%. And that's assuming your outcomes are limited to those two, and that the probability is accurate. With $p = 44.4\%$, Kelly gives zero, so only 16 points of error take the size from 3.5 times to nothing.

\iqwrong{Giving the Kelly number as your size. It needs a true probability, and it's your estimate of the probability that the interviewer is testing.}
\end{solution}

\begin{solution}{iq:iv2:your-strategy-and-your-project:13}
About 14\% of flagged cases are frauds. In 1\,000 cases, I'd expect 10 frauds, of which 8 are caught, and about 50 good cases flagged: 58 flags, of which 8 are right.

The precision is $0.008/(0.008 + 0.0495) = 13.9\%$.

So an AUC of 0.95 says little about use. Show a precision and recall curve at the base rate, and the cost of each kind of error. Then the threshold choice is a business decision, and you can defend it.

\iqwrong{Leading with AUC or accuracy on rare events. Always saying ``no fraud'' gets 99\% accuracy.}
\end{solution}

\begin{solution}{iq:iv2:your-strategy-and-your-project:14}
About 2 points of the 80\% are selection if the submissions were independent, and about 0.64 at a correlation of 0.9. To separate you from a rival 0.3 points behind, you'd need about 137\,000 rows with independent errors, and about 27\,000 with errors correlated 0.8.

\textbf{Selection.} One public score has a standard error of $\sqrt{0.8 \times 0.2/2000} = 0.894$ points. The best of 50 independent draws has an expected value of 2.25 standard errors. So $2.25 \times 0.894 = 2.0$ points.

\textbf{Correlated.} With equal correlation $\rho$, only the independent part spreads, so the factor is $\sqrt{1 - \rho} = 0.316$, which gives 0.64 points.

\textbf{Separation.} With independent errors, the gap has a standard error of $\sqrt{2 \times 0.16/N}$. For $0.003$ to be 1.96 of these, $N = 2 \times 0.16 \times (1.96/0.003)^2 = 136\,600$. On the same rows, with errors correlated across models, the variance of the gap falls with the correlation: at 0.8, it is $2 \times 0.16 \times 0.2/N$, which needs about 27\,300 rows.

So a rank of 10th of 2\,000 is about the workflow, and the honest way to report it is with the final score on data you couldn't see.

\iqwrong{Quoting your public score, or a place, as your skill. Ask what the private set said, and how many others were within one standard error.}
\end{solution}

\begin{solution}{iq:iv2:your-strategy-and-your-project:15}
You'd expect about 19\%, and only about 7\% of paths do as well as 12\%. So it isn't good news: either the sample was gentle, or the record was chosen, and you shouldn't size for 12\%. On a normal day, you'd expect to be down about 5 to 6\% from the high.

\textbf{Simulation.} With a daily normal walk at those numbers, the 5-year maximum has a mean of about 19\%, a median of about 18\%, and a 95th percentile of about 31\%. A drawdown past 30\% has about a 6\% chance even if the strategy is exactly as good as you say.

\textbf{Stationary formula.} For a Brownian motion with drift $\mu$ and volatility $\sigma$, the distance from the running high has an exponential distribution in the long run, with mean $\sigma^2/(2\mu) = 0.0225/0.4 = 5.6\%$. This also says that more drift cuts drawdowns fast, and that they grow with the square of the volatility.

So when you set a stop-out level, you should use a number you'd expect, not your own lucky one. On that simulation, a 12\% limit shuts down a healthy strategy about 93 times in 100 within five years.

\iqwrong{Quoting your sample's drawdown as your risk. It's one draw from a wide distribution, and on the low side here.}
\iqfollow{What if the true return is half?}{The expected drawdown rises to about 26\%, since the drift holds it back less. Redo the simulation with the new drift before you set a limit.}
\end{solution}

\begin{solution}{iq:iv2:your-strategy-and-your-project:16}
The test shows a mean of about 1 microsecond and a 99th percentile of 1 microsecond. A stream on its own schedule sees a mean of about 56 microseconds, a 99th percentile of about 1 millisecond, and a 99.9th of about 9 milliseconds. You should have measured from each message's intended send time.

\textbf{What you measured.} In a second, you send about 990\,000 messages at 1 microsecond, and one takes 10 milliseconds. So one slow sample is hidden among 990\,000 good ones, and the mean is 1.01 microseconds.

\textbf{The stream.} A message that arrives $t$ milliseconds into the stall waits for the rest of the stall and for the $0.1t$ of work ahead of it: $10 - 0.9t$ milliseconds. The queue keeps draining at that rate after the stall ends, so the wait falls in a straight line from 10 milliseconds to zero over $10/0.9 = 11.1$ milliseconds. The mean is $\tfrac12 \times 10 \times 11.1/1000 = 56$ microseconds. For the tail, the fraction of the second with a wait above $w$ is $(10 - w)/0.9$ thousandths. Setting that to 1\% gives $w = 1$ millisecond, and to 0.1\% gives $w = 9.1$ milliseconds.

\textbf{What to do.} The trouble is called coordinated omission: the test stops sending just when the system is bad. So use an open-loop generator with a fixed schedule, measure from when each message should have gone, and report the tail with the rate. Then find the stall.

\iqwrong{Quoting a mean, or saying that one slow message in a second doesn't matter. In a trading system, the stall is when the prices are moving, and the messages all arrive at that moment.}
\end{solution}

\begin{solution}{iq:iv2:your-strategy-and-your-project:17}
It's about 0.95 with those days and about 0.39 without. A normal distribution would give about 14 points from the top 5 days, not 40. So the result depends on a few events, and the question is whether you can explain them.

\textbf{Sharpes.} With them, the mean is 6 basis points a day on a volatility of 100, and $0.06 \times \sqrt{252} = 0.95$. Without them, the mean is $20/995 = 2.0$ basis points, and the volatility falls too, because each of those days is 8 points: the variance loses $5 \times 64$ of about 1\,004 square points, so the volatility is about 83 basis points. That gives $(2.01/82.9) \times \sqrt{252} = 0.385$.

\textbf{Normal.} Removing the top 5 from any sample lowers its mean, so you must compare with the normal case. The top 0.5\% of a normal has a mean of $\varphi(2.576)/0.005 = 2.89$ standard deviations. So $5 \times 2.89 = 14.5$ points. The exact finite-sample number is about 14.3. A symmetric normal strategy would still fall by 14 points without them, and this one loses 40.

\textbf{What it tells you.} Each of the top days is about 8 standard deviations. That's either a real source of edge you can name (such as a crisis the strategy is built to profit from), or a data error, or luck. Check the five dates first. If they're real, then show that the same kind of event should come again, and how often. If you remove the top days, also remove the bottom ones, or the comparison is unfair.

\iqwrong{Saying the result is robust since it's positive without the best days. It's the size of the drop against a normal one that says how dependent it is.}
\end{solution}

\begin{solution}{iq:iv2:your-strategy-and-your-project:18}
It's about 72\%. At a 5\% level, it's about 34\%, and at a 0.1\% level, about 96\%.

Of 1\,000 ideas, 50 are true, and half of these are found: 25. Of the 950 false, 1\% pass: 9.5. So $25/34.5 = 72\%$.

At 5\%: $25/(25 + 47.5) = 34\%$. At 0.1\%: $25/(25 + 0.95) = 96\%$, assuming the same power, which isn't realistic, as a stricter level loses power.

\textbf{What I say.} ``It's significant at 1\%, which on my guess at the field's base rate gives a true result about 7 times in 10. Not 99\%. It's also one test, and I'd want it out of sample.''

\iqwrong{Saying that a 1\% $p$-value means a 99\% chance that it's real. That confuses the chance of the data given no effect with the chance of no effect given the data.}
\end{solution}

\begin{solution}{iq:iv2:your-strategy-and-your-project:19}
It's about 1.86 if you assume independence, and about 1.37 with the autocorrelation. In one sentence: my model's average loss was lower, but with 500 days I can't rule out that it's noise.

$t = 0.05/0.6 \times \sqrt{500} = 1.86$. For an AR(1) with coefficient $\rho$, the long-run variance is multiplied by $(1 + \rho)/(1 - \rho) = 1.857$, so the $t$ is divided by 1.363.

This is the Diebold-Mariano test: a $t$-test on the loss differences using a long-run variance. And the two forecasts need to be compared fairly, on the same dates, with neither tuned on the test period.

\iqwrong{Saying ``it beat the benchmark'' on the mean alone, or using the plain standard error on data whose errors are persistent.}
\end{solution}

\begin{solution}{iq:iv2:your-strategy-and-your-project:20}
(a) 3.0. (b) A return of about 14.2\% and a Sharpe of about 1.19. (c) Alpha of 14.4\%, and an information ratio of about 1.31. (d) Not by itself: 1.4 is only about 1.4 standard errors. (e) About 37. (f) I believe a Sharpe of about 1, not 1.5, and I'd run it small until the fills and second-half decay are explained.

\textbf{(a)} The Sharpe is $18/12 = 1.5$, so $t = 1.5\sqrt{4} = 3.0$.

\textbf{(b)} Cost is $0.5 \times 3 \times 252 = 378$ basis points a year. Doubling it takes $18 - 3.78 = 14.22\%$, and the Sharpe is 1.19.

\textbf{(c)} $18 - 0.3 \times 12 = 14.4$. Residual variance is $0.0144 - 0.09 \times 0.0256 = 0.0121$, so its root is 11\%, and the ratio is $14.4/11 = 1.31$.

\textbf{(d)} Each half is 2 years, so each has a standard error of $1/\sqrt{2}$. Their difference has a standard error of 1. So $2.2 - 0.8 = 1.4$ is 1.4 standard errors, with a two-sided $p$ of 16\%. It's not evidence of decay, nor is it evidence against it. Look at the signal and the costs, which are more telling than the P\&L.

\textbf{(e)} One-sided, $\Phi(-3) = 0.135\%$. For the chance of at least one false pass among $N$ to reach 5\%: $N = \ln 0.95/\ln(1 - 0.00135) = 38$, or $0.05/0.00135 = 37$ by the simple rule. So I must be able to say that I tried far fewer than that, and I must count the correlated variants as well.

\textbf{(f)} Say something like this. ``Three things pull the number down: the second half, a realistic cost, and my search. I don't have a test that separates them in four years. So I treat this as a Sharpe of about 1, and the net of costs is what I'd check first live, because cost per dollar can be measured in weeks. I would size it for that, and not for 1.5. I'd stop if the signal's measured information coefficient fell, or if live costs went past double the model. Those are checkable. A bad quarter isn't.''

\iqwrong{Answering each question well and never reaching a conclusion. The last question is the one they care about, and you should give a number and a plan.}
\iqfollow{You're telling me that the Sharpe of 1 is a judgement. How did you get it?}{Be open that it's judgement: the 1.5 reduced for cost realism, and a selection haircut that depends on the search. Show the range, and say what would move you.}
\end{solution}
